\section{Hyperparameter defaults}

本节整理 FFN ratio、head dim、aspect ratio、vocab size、dropout、weight decay 的现代默认值和例外。


\begin{knowledgebox}{超参数默认值的读法}
本节所有比例都应读作“保守起点”，不是“唯一正确”。若模型规模、数据、硬件、上下文长度或推理目标不同，就需要用小规模实验验证这些默认值是否仍适用。
\end{knowledgebox}


架构模块确定之后，仍有一组“没有新模块、但会决定模型形状”的问题：FFN 应放大多少倍、head 数与 head dimension 如何配合、词表多大、模型更深还是更宽、是否使用 dropout。Slide 36 的作用是把这些问题显式化，避免默认值悄悄藏在配置文件里，最后被误认为与质量或成本无关。

\begin{figure}[H]
\centering
\includegraphics[width=0.88\textwidth]{slide-035.jpg}
\caption{Slide 36：Transformer 的关键 shape 与 regularization 超参数问题。}
\figsource{CS336 Spring 2026 Lecture 3 官方幻灯片。}
\end{figure}

\begin{knowledgebox}{超参数其实是资源分配规则}
FFN ratio 决定多少参数和 FLOPs 放在 token-wise MLP；head shape 决定 attention projection 与 kernel shape；vocab size 改变 embedding/output matrix 和 token 序列长度；depth/width 影响并行、latency 与优化路径。它们不是互相独立的旋钮，任何一个变化都可能重分配训练和推理成本。
\end{knowledgebox}


\begin{figure}[H]
\centering
\includegraphics[width=0.88\textwidth]{slide-036.jpg}
\caption{Slide 37: FFN ratio consensus}
\figsource{CS336 Spring 2026 Lecture 3 官方幻灯片。}
\end{figure}

\noindent\textbf{展开说明：}本页展示 FFN hidden dim 与 model dim 的常见比例，传统 dense FFN 常用约 4 倍。



普通两层 FFN 常取 $d_{ff}\approx4d_{model}$；但 gated FFN 多了一条输入投影，若仍保持 4 倍中间维度，参数量与 FLOPs 会显著增加。Slide 38 因此比较公开模型实际使用的 $d_{ff}/d_{model}$：LLaMA、Qwen、DeepSeek、Yi 等大量 GLU 模型落在约 2.5--3.5 的范围，接近按等预算推导的 $8/3$，而 PaLM、Mistral 与 LLaMA 2 稍大。

\begin{figure}[H]
\centering
\includegraphics[width=0.88\textwidth]{slide-037.jpg}
\caption{Slide 38：GLU 模型的 FFN multiplier 与等参数预算下的 $8/3$ 经验值。}
\figsource{CS336 Spring 2026 Lecture 3 官方幻灯片。}
\end{figure}

\begin{importantbox}{为什么是约 $8/3$，而不是固定 $4$}
普通 FFN 主要有两张矩阵，参数规模约为 $2d_{model}d_{ff}$；GLU 有 value、gate 和 output 三张矩阵，约为 $3d_{model}d_{ff}^{GLU}$。令二者接近，得到 $d_{ff}^{GLU}\approx\frac{2}{3}d_{ff}$；若普通 FFN 用 $4d_{model}$，则 GLU 约为 $\frac{8}{3}d_{model}$。实际模型还会为了 tensor-core 对齐取整。
\end{importantbox}

默认范围并不意味着所有模型都保守。Slide 39 用 T5 11B 的 $d_{ff}=65{,}536$ 与 $d_{model}=1{,}024$ 展示一个 64 倍极端 aspect，随后列出 Gemma 系列等较新的更大 multiplier。教学重点不是复制 T5 数字，而是认识到超参数表中的例外可能来自时代、训练目标或系统实现；遇到离群点时应追查完整 recipe，而不是单独迁移一个 ratio。

\begin{figure}[H]
\centering
\includegraphics[width=0.88\textwidth]{slide-038.jpg}
\caption{Slide 39：T5 的极端 FFN multiplier 与其他近期例外。}
\figsource{CS336 Spring 2026 Lecture 3 官方幻灯片；Raffel et al. 2020。}
\end{figure}

\begin{warningbox}{离群配置不能脱离参数账本}
一个超大 FFN multiplier 会把参数和计算大量移入 MLP，同时改变 attention/MLP 的比例、激活内存和并行方式。若只把 multiplier 抄到另一模型，却不匹配总参数、训练 token、optimizer 与硬件 shape，得到的不是对原 recipe 的复现。
\end{warningbox}


\begin{figure}[H]
\centering
\includegraphics[width=0.88\textwidth]{slide-039.jpg}
\caption{Slide 40: Why multiplier range}
\figsource{CS336 Spring 2026 Lecture 3 官方幻灯片。}
\end{figure}

\noindent\textbf{展开说明：}本页展示 ff\_dim multiplier 在一定盆地内都可行，超参选择不是尖锐单点，而是宽容区间。

\begin{knowledgebox}{读图：Slide 40 应该怎么看}
这页不是装饰性截图，而是本节论点的证据页。先看图中模型/论文/曲线或表格的比较对象，再看它们支持的趋势：哪些设计已经形成共识，哪些只是局部实验结果，哪些主要由运行时或推理成本推动。读这类页时要区分 quality evidence、runtime evidence 和 stability evidence，不能把一种证据直接替换成另一种。
\end{knowledgebox}


\begin{figure}[H]
\centering
\includegraphics[width=0.88\textwidth]{slide-040.jpg}
\caption{Slide 41: Model-dim lesson}
\figsource{CS336 Spring 2026 Lecture 3 官方幻灯片。}
\end{figure}

\noindent\textbf{展开说明：}本页总结：保守默认值好用，但应通过小规模实验验证，不要把公式当法律。


\begin{figure}[H]
\centering
\includegraphics[width=0.88\textwidth]{slide-041.jpg}
\caption{Slide 42: Head dim consensus}
\figsource{CS336 Spring 2026 Lecture 3 官方幻灯片。}
\end{figure}

\noindent\textbf{展开说明：}本页讨论 head\_dim × num\_heads 与 model\_dim 的关系，现代模型通常让总 head dim 接近 model dim。



抽象关系之后，Slide 43 列出 GPT-3、T5、PaLM、LLaMA 2 与 Qwen 3.5 等实例。多数模型让 $n_{heads}\times d_{head}$ 与 $d_{model}$ 同阶，但 T5 等 Google 模型出现明显偏离。读表时应同时看 head dimension 是否便于 kernel、总 Q/K/V projection 是否扩大，以及后续是否采用 GQA/MQA；单独看 head 数没有可比性。

\begin{figure}[H]
\centering
\includegraphics[width=0.88\textwidth]{slide-042.jpg}
\caption{Slide 43：不同模型的 head 数、head dimension、model dimension 与比例。}
\figsource{CS336 Spring 2026 Lecture 3 官方幻灯片。}
\end{figure}

\begin{knowledgebox}{读表：head count 不是容量的独立指标}
增加 head 数若同时减小每个 head 的维度，attention projection 的总宽度可能不变；反之，总 head dimension 大于 model dimension 会扩张投影成本。比较模型时至少记录 $d_{model}$、$d_{head}$、query head 数、KV head 数和 projection ratio，否则无法判断质量与 KV cache 成本。
\end{knowledgebox}

接下来转向整体 aspect ratio，也就是模型更深还是更宽。Slide 44 用 $d_{model}/n_{layer}$ 汇总公开模型，显示许多主流 recipe 落在相对窄的范围，但 BLOOM、T5 与 Gemma 等仍有明显差异。这个比值只是描述性统计：相同值不代表相同参数量，embedding、FFN ratio 和 attention projection 仍会改变总预算。

\begin{figure}[H]
\centering
\includegraphics[width=0.88\textwidth]{slide-043.jpg}
\caption{Slide 44：公开模型的 depth-width aspect ratio。}
\figsource{CS336 Spring 2026 Lecture 3 官方幻灯片。}
\end{figure}

\begin{knowledgebox}{Aspect ratio 的三重约束}
更深会增加串行层数、pipeline 边界和逐 token latency，但每层矩阵较小；更宽可用更大 GEMM 提高单层效率，却增加 activation、通信量和单层参数。优化上，深度还改变 residual path 与梯度传播。因此最佳 ratio 必须同时满足统计效率、并行拓扑和服务延迟。
\end{knowledgebox}

Slide 45 专门把系统代价画出来：极深模型更难做低延迟推理，也更难在 pipeline parallelism 中保持均衡；每个 token 必须依次穿过更多层，无法像 batch 维度那样完全并行。课程由此把 aspect ratio 从“神秘架构偏好”还原为可测的 latency/parallelism tradeoff。

\begin{figure}[H]
\centering
\includegraphics[width=0.88\textwidth]{slide-044.jpg}
\caption{Slide 45：极深模型的并行化与推理延迟代价。}
\figsource{CS336 Spring 2026 Lecture 3 官方幻灯片；Tay et al. 2021。}
\end{figure}

\begin{warningbox}{只匹配参数量仍可能比较不公平}
两个参数量相同的模型若 depth-width 不同，训练 step time、pipeline bubble、activation checkpointing 成本和 decode latency 都可能不同。架构消融应至少同时报告参数量、训练 FLOPs、tokens/s、显存峰值与服务 latency，不能只报告 validation loss。
\end{warningbox}


\begin{figure}[H]
\centering
\includegraphics[width=0.88\textwidth]{slide-045.jpg}
\caption{Slide 46: Aspect ratio evidence}
\figsource{CS336 Spring 2026 Lecture 3 官方幻灯片。}
\end{figure}

\noindent\textbf{展开说明：}本页引用 Kaplan/Tay 等证据说明 depth/width scaling 有经验趋势，但仍依赖任务和规模。

\begin{knowledgebox}{读图：Slide 46 应该怎么看}
这页不是装饰性截图，而是本节论点的证据页。先看图中模型/论文/曲线或表格的比较对象，再看它们支持的趋势：哪些设计已经形成共识，哪些只是局部实验结果，哪些主要由运行时或推理成本推动。读这类页时要区分 quality evidence、runtime evidence 和 stability evidence，不能把一种证据直接替换成另一种。
\end{knowledgebox}


\begin{figure}[H]
\centering
\includegraphics[width=0.88\textwidth]{slide-046.jpg}
\caption{Slide 47: Vocabulary size}
\figsource{CS336 Spring 2026 Lecture 3 官方幻灯片。}
\end{figure}

\noindent\textbf{展开说明：}本页对比单语、多语、生产系统的 vocab size，显示 tokenizer/词表也是架构超参。


\begin{figure}[H]
\centering
\includegraphics[width=0.88\textwidth]{slide-047.jpg}
\caption{Slide 48: Dropout regularization}
\figsource{CS336 Spring 2026 Lecture 3 官方幻灯片。}
\end{figure}

\noindent\textbf{展开说明：}本页询问预训练是否需要 dropout。大型 LM 预训练常少用 dropout，因为数据量巨大且目标不是小数据过拟合。



Slide 49 把历史实践放进同一张表：较早的 Transformer、GPT-2/3 与 OPT 常在预训练使用 dropout，而 T5 v1.1、PaLM、LLaMA 等较新大模型更常设为零；weight decay 则被许多 recipe 保留。这里存在报告偏差——论文不写 dropout 往往意味着开源配置中没有，但不能据此断言所有闭源模型都相同。

\begin{figure}[H]
\centering
\includegraphics[width=0.88\textwidth]{slide-048.jpg}
\caption{Slide 49：公开模型在预训练阶段使用 dropout 与 weight decay 的实践。}
\figsource{CS336 Spring 2026 Lecture 3 官方幻灯片。}
\end{figure}

\begin{knowledgebox}{读表：regularization 选择随数据规模变化}
大规模预训练通常不缺随机性，数据量也远大于小型监督任务，因此 dropout 的收益下降且可能降低吞吐；weight decay 则直接约束参数尺度，并与 optimizer 动力学相连。这个趋势不应迁移到小数据 fine-tuning：数据变少、训练轮数增加后，dropout 或其他 regularization 可能重新有价值。
\end{knowledgebox}


\begin{figure}[H]
\centering
\includegraphics[width=0.88\textwidth]{slide-049.jpg}
\caption{Slide 50: Why weight decay LLMs}
\figsource{CS336 Spring 2026 Lecture 3 官方幻灯片。}
\end{figure}

\noindent\textbf{展开说明：}本页引用 Andriushchenko 等工作，提示 weight decay 可能影响 logit/表示尺度和泛化。


\begin{figure}[H]
\centering
\includegraphics[width=0.88\textwidth]{slide-050.jpg}
\caption{Slide 51: Hyperparameter summary}
\figsource{CS336 Spring 2026 Lecture 3 官方幻灯片。}
\end{figure}

\noindent\textbf{展开说明：}本页总结 FFN ratio、head ratio、aspect ratio、vocab、dropout、weight decay 的现代保守默认。

\subsection{本章小结}
本节的关键不是记住所有模型名，而是把公开模型的设计选择转化成可迁移判断：共识默认值可以作为起点，例外需要知道原因，涉及系统和推理成本的选择必须结合资源账本讨论。

